% =====================================================================
%  Section II. Quark Binding Mechanism
% =====================================================================
\section{QUARK BINDING MECHANISM}\label{sec:binding}

The binding mechanism will be explained in this section using the
Feynman path-integral picture.  The path-integral framework will be used
in an intuitive rather than a formal way.  Consider the current-current
propagator
\begin{equation}
D_{\mu\nu}(x) = \ev{\mathrm{T}\,J_{\mu}(x)J_{\nu}(0)}\,,
\label{eq:2.1}
\end{equation}
whose Fourier transform determines the $e^{+}e^{-}$ annihilation cross
section into hadrons.  Assume that the currents $J_{\mu}(x)$ are built
from quark fields as in the quark-parton model.  Assume that the quarks
interact through a single gauge field.  (The restriction to one field is
only for simplicity.)  In the Feynman path-integral picture the
propagator $D_{\mu\nu}(x)$ is given by a weighted average over all
possible classical quark paths and all possible classical values of the
gauge field.  The currents $J_{\mu}(x)$ and $J_{\nu}(0)$ are thought of
as producing a quark pair at the origin which later annihilate at the
point $x$:  One has to sum over all paths joining the points $0$ and $x$
for each of the pair of quarks.

An example of paths for the quark and antiquark are shown in
Fig.~\ref{fig:3}.  The vacuum can also emit and absorb quark pairs; this
leads to further closed loops as illustrated in Fig.~\ref{fig:4}.  All
possible loops must be summed over too.  It is also possible to have
independent loops for the points $0$ and $x$ (Fig.~\ref{fig:5}), but
this possibility will not be important here.

\begin{figure}[t]
  \centering
  \includegraphics[width=0.52\textwidth]{fig3}
  \refstepcounter{figure}%
  \caption*{Fig. 3. An example of quark ($q$) and antiquark ($\bar q$)
    paths connecting the points $0$ and $x$.}
  \label{fig:3}
\end{figure}

\begin{figure}[t]
  \centering
  \includegraphics[width=0.5\textwidth]{fig4}
  \refstepcounter{figure}%
  \caption*{Fig. 4. Example of current loop (as in Fig. 3) with extra
    vacuum loop.}
  \label{fig:4}
\end{figure}

The weight associated with a given quark path or set of paths includes a
factor of $\exp[ig\oint A_{\mu}(x)\,ds^{\mu}]$, where $A_{\mu}(x)$ is
the gauge field.  Here $\oint ds^{\mu}$ is a line integral or a sum over
line integrals for each of the quark-antiquark loops, including the loop
joining the points $0$ and $x$.  The constant $g$ is the coupling
constant of the gauge theory.  There are further weight factors
independent of $A$.  Finally, independently of the quark paths there is
another weight factor, namely, the exponential of the free action for
the gauge field.  The combined weight factor is then averaged over all
quark paths and all gauge fields $A_{\mu}(x)$ to give the current-current
propagator.

\begin{figure}[t]
  \centering
  \includegraphics[width=0.42\textwidth]{fig5}
  \refstepcounter{figure}%
  \caption*{Fig. 5. Example of separate quark loops for the points $0$
    and $x$. (Integration over the gauge field produces gauge
    propagators which connect these loops.)}
  \label{fig:5}
\end{figure}

In order that quarks exist as separate final-state particles it must be
possible to have quark-antiquark loops with well-separated quark and
antiquark lines, at least when $x$ and $0$ are far apart.  This is
illustrated in Fig.~\ref{fig:6}(a).  If the quark and antiquark paths
are unlikely to separate beyond a fixed size, say $10^{-13}$~cm [see
Fig.~\ref{fig:6}(b)], then clearly no detector will see a quark or
antiquark in isolation.

It is assumed in this discussion that vacuum loops are not important.
If vacuum loops are important enough then space will be filled with a
high density of vacuum-produced quark-antiquark pairs.  In particular,
there will be many quark-antiquark pairs inside a detector of
macroscopic size.  The question then is whether there can be an excess
of quarks over antiquarks, or vice versa, in a region of macroscopic
size.  This is a more difficult question to answer and will not be
discussed in this paper.

\begin{figure}[t]
  \centering
  \includegraphics[width=0.55\textwidth]{fig6}
  \refstepcounter{figure}%
  \caption*{Fig. 6. (a) Loop with well-separated quark and antiquark.
    (b) Loop with small separation between quark and antiquark.}
  \label{fig:6}
\end{figure}

Note that $x$ must be large:  If $x$ is small there is little likelihood
of finding a large size loop.  This may seem a bit peculiar:  One expects
quarks to appear in the final state of $e^{+}e^{-}$ annihilation only at
large virtual-photon momentum $q$ if they appear at all, and large $q$
means small $x$, not large $x$.  The answer to this paradox has two
parts.  First, the important paths in the Feynman path integral bear no
detailed relation to possible physical final states (the paths are paths
of bare particles, not physical particles).  Secondly, large $x$ does
not necessarily mean small $q$.  In fact the study of whether
well-separated quark-antiquark paths exist for large $x$ is really a
search for a quark-antiquark threshold in $e^{+}e^{-}$ annihilation,
which would contribute a term $\sim \exp(2imx)/x^{p}$ to the
current-current propagator for large $x$, where $m$ is the quark mass.
(Here $\sim$ means up to a power of $x$.${}^{9}$)  Such a term
corresponds in momentum space to the singularity at the threshold
$q^{2} = (2m)^{2}$.

Suppose the gauge-field averaging is performed before the quark-paths
averaging.  Then one determines the average over all gauge fields of the
weight factor $\exp[ig\oint A_{\mu}(y)\,ds^{\mu}]$ weighted further with
the exponential of the free gauge-field action.  For an Abelian gauge
theory this average can be computed explicitly:  It is
\[
  \exp\!\Big[-g^{2}\oint\!\!\oint ds^{\mu}\,ds'^{\nu}\,D_{\mu\nu}(y-y')\Big],
\]
where $D_{\mu\nu}(y-y')$ is the free gauge-field propagator.

The quark binding mechanism can be seen by comparing the above
expression for one space dimension and three space dimensions.  In three
space dimensions this calculation gives no binding (the binding occurs
only with a modified gauge-field action:  see Sec.~IV), while there is
binding in one space dimension.  In three space dimensions
$D_{\mu\nu}(y-y')$ behaves as $(y-y')^{-2}$.  In consequence large
values of $(y'-y)$ are negligible in the double line integral.  Hence,
the double integral is proportional to $P$, where $P$ is the length of
the loop.  Unfortunately the integral is divergent at $y'=y$; a cutoff
is needed for the integral to make sense.  Since a cutoff will be
introduced anyway in this paper, this is not a major concern.  For
simple loops, the perimeter $P$ is roughly of order $(x^{2})^{1/2}$
(ignoring the case that $x$ is close to the light cone).  Thus one has
an exponential of the type one expects when free quarks are present.

In one space dimension, $D_{\mu\nu}(y-y')$ behaves as
$\ln[(y-y')^{2}]$ for $y'-y$ large, and $y'$ and $y$ can freely range
separately over the loop.  In this case the double integral is
proportional to $P^{2}$.  Now the gauge-field average behaves as
$e^{-cP^{2}}$ where $c$ is a constant.  In this case the contribution of
large loops is heavily suppressed and there are no free quarks.  [The
case of nearby quark-antiquark pairs as in Fig.~\ref{fig:6}(b) is
special --- in this case large $y-y'$ is unimportant due to cancellation
between the quark path and the nearby and oppositely directed antiquark
path.  In this case the double integral behaves as $P$, not $P^{2}$, but
in this case there are no isolated quarks.]

In the strongly coupled lattice gauge theory described in later
sections, the gauge-field average of $\exp[ig\oint A_{\mu}(x)\,ds^{\mu}]$
behaves as $\exp(-cA)$, where $A$ is the enclosed area of the loop.
This heavily suppresses large loops, such as in Fig.~\ref{fig:6}(a),
where $A$ is of order $P^{2}$.  One can think of one factor $P$ as being
roughly $(x^{2})^{1/2}$, the other $P$ as being analogous to a mass
multiplying $(x^{2})^{1/2}$.  Since $P \to \infty$ as $x \to \infty$,
the quark-antiquark threshold is at infinite mass.

\begin{figure}[t]
  \centering
  \includegraphics[width=0.42\textwidth]{fig7}
  \refstepcounter{figure}%
  \caption*{Fig. 7. Quark-antiquark loop with nearby vacuum loop.}
  \label{fig:7}
\end{figure}

In all these calculations one can have a large loop if there is a nearby
vacuum loop (Fig.~\ref{fig:7}).  In this case one always gets
$e^{-cP}$ behavior.  For example, in the strong-coupling case the
relevant enclosed area is the area between the two loops which is
proportional to the perimeter $P$ provided the separation of the two
loops is fixed independently of $P$.  This is in accord with Schwinger's
picture.  While an isolated well-separated loop may be suppressed (due
to $P^{2}$ or $A$ dependence in the exponential) a loop closely shielded
by a vacuum-polarization loop is always unsuppressed.

The binding mechanism proposed here is soft:  The exponential damping is
associated with large size loops having large areas.  The behavior of
small loops is irrelevant to the binding mechanism.  Also for small
loops both their area and perimeter are small and neither is of great
importance in an exponential.

The mechanism discussed here works equally well for Dirac quarks or
scalar quarks.  This is in contrast to the Higgs mechanism which can
wipe out the charge of scalar particles only.
